\chapter{Теория}

\section{Разделение на 5 сетей}

\begin{figure}[H] 
    \centering
    \includegraphics[width=0.75\textwidth]{start.png}
    \caption{Начальные данные}
\end{figure}

Для того, чтобы можно было образовывать новые сети, нужно занять в сетевую часть биты из узловой части. Тогда сеть можно будет
разделить на $2^n$ сетей, где n - кол-во занятых бит. Необходимо занять 3 бита, тогда получим следующие данные:

\begin{figure}[H] 
    \centering
    \includegraphics[width=0.75\textwidth]{start2.png}
    \caption{Перераспределение бит из сетевой и узловой части}
\end{figure}

\subsection{Первая сеть}

Чтобы определить диапазон адресов, нужно заполнить все узловые биты единицами и отнять 2. Это будет правая границы сети. \\

Адрес сети: 10.1.0.0\\
Диапазон адресов: 10.1.0.1 - 10.1.31.254

\subsection{Вторая сеть}

Адрес сети: 10.1.32.0\\
Диапазон адресов: 10.1.32.1 - 10.1.63.254

И т.д

\endinput